\documentclass[11pt,a4paper,chinese,fontset=overleaf]{ctexart} % 文档类型为ctexart，基础字号 11 磅，纸张尺寸 A4，支持中文
% \documentclass[11pt,a4paper,chinese,fontset=windows]{ctexart}
% 加载各种宏包
\usepackage{array,lastpage,amsmath,amssymb,mathtools,enumitem,graphicx,multirow,tocbibind,longtable,makecell,varwidth,titlesec,bm,booktabs,comment,qtree,listings,tikz,flowchart,algorithm,algorithmic}
\usepackage{olymp}
\usepackage{fancyhdr}
 
% \setCJKmainfont[BoldFont=Heiti SC Medium]{Songti SC Light}
% \setCJKmainfont[BoldFont=SimHei]{SimSun}

% 页眉页脚设置
\pagestyle{fancy}
\fancyhf{}
\chead{
    \sffamily\normalsize {{CONTEST_TITLE}}（{{CONTEST_STAGE}}）\\
    \sffamily\normalsize {{CONTEST_DATE_HEADER}}
} % 页眉：比赛名称与日期
\renewcommand{\headrulewidth}{.5pt} % 页眉线宽
\renewcommand{\footrulewidth}{.5pt} % 页脚线宽
\setlength{\headsep}{20pt} % 页眉与正文间距
\cfoot{\sffamily \normalsize Page \thepage \ of \pageref{LastPage}} % 页脚居中显示页码  包含第一页

% 居中显示 verbatim 内容（如输入输出样例）
\usepackage{varwidth}
\newenvironment{centerverbatim}{\par\centering\varwidth{\linewidth}\verbatim}{\endverbatim\endvarwidth\par}

\usepackage[thinlines]{easytable}
% 代码高亮配置：代码自动换行、禁用边框、使用等宽字体
\lstset{escapeinside=``, breaklines=true, frame=none, extendedchars=false, basicstyle=\ttfamily, showstringspaces=false}
% 优化中文排版
\setlength{\parindent}{2em} % 首行缩进 2 字符
\setlength{\parskip}{2.0ex} % 段间距 2 倍行高
\linespread{1.15} % 行距 1.15 倍

% 数学符号简写
\newcommand*{\mat}[1]{\boldsymbol{\mathrm{#1}}} % 粗体矩阵
\renewcommand*{\vec}[1]{\boldsymbol{\mathrm{#1}}} % 粗体向量

% 绘图工具配置
\usepackage{tikz,flowchart}
\usetikzlibrary{shapes,shapes.geometric,arrows,automata,positioning}
\usepackage{pgfplots}

% Keep the supplied main.tex typography, offsets, header/footer and fontset.
% Administrator overrides are optional; the archive preset has none.

